% Compile with: uplatex file.tex && dvipdfmx file.dvi
%           or: latexmk -uplatex file.tex
\documentclass[dvipdfmx]{article}
\usepackage{amsmath, amssymb}
\usepackage{tikz}

\title{scenario\_hubbard\_dimer\_operator\_algebra}
\date{\today}

\begin{document}
\maketitle

Declare the electron annihilation operator, the hopping $w$, the on-site repulsion $U$ and the time $t$.

define c as an electron annihilation operator

\[
  c
\]

define w as a scalar parameter

\[
  w
\]

define U as a scalar parameter

\[
  U
\]

define t as a scalar parameter

\[
  t
\]

\textbf{The two-site Hubbard model, written in the textbook sum notation.} The spin sum $\sum_\sigma$ carries no range, so it goes to QuantumAlgebra's formal sum as written; the site sum $\sum_{i=1}^{2}$ has a concrete range and is therefore \textbf{expanded into terms} before it reaches the backend (QA's $\sum$ has nowhere to put a range — which is no reason to drop the range silently and compute something other than what was written).

$H = -w \sum_{\sigma} (c^\dagger_{1\sigma} c_{2\sigma} + c^\dagger_{2\sigma} c_{1\sigma}) + U \sum_{i=1}^{2} c^\dagger_{i\uparrow} c_{i\uparrow} c^\dagger_{i\downarrow} c_{i\downarrow}$

\begin{equation}
  H = -w \sum_{\sigma} (c^\dagger_{1\sigma} c_{2\sigma} + c^\dagger_{2\sigma} c_{1\sigma}) + U \sum_{i=1}^{2} c^\dagger_{i\uparrow} c_{i\uparrow} c^\dagger_{i\downarrow} c_{i\downarrow}
  \label{eq:E1}
\end{equation}

Ask for the underlying anticommutators \textbf{in the bracket notation itself}: 1 for the same site and spin, 0 for a different site, and \textbf{0 for a different spin} — that last one stays as $\delta_{\uparrow\downarrow}$ unless the arrow is read as a value.

$\{c_{1\uparrow}, c^\dagger_{1\uparrow}\}$

\begin{equation}
  1
  \label{eq:E2}
\end{equation}

$\{c_{1\uparrow}, c^\dagger_{2\uparrow}\}$

\begin{equation}
  0
  \label{eq:E3}
\end{equation}

$\{c_{1\uparrow}, c^\dagger_{1\downarrow}\}$

\begin{equation}
  0
  \label{eq:E4}
\end{equation}

\textbf{Particle-number conservation.} Write the total number $\hat N=\sum_{i\sigma}\hat c^\dagger_{i\sigma}\hat c_{i\sigma}$ and compute $[\hat H,\hat N]$. The $\sum_\sigma$ expands over the projections of the particle that was DECLARED (an electron, $m=\pm1/2$). Neither hopping nor $U$ changes $\hat N$, so it should be 0.

$N = \sum_{i=1}^{2} \sum_{\sigma} c^\dagger_{i\sigma} c_{i\sigma}$

\begin{equation}
  N = \sum_{i=1}^{2} \sum_{\sigma} c^\dagger_{i\sigma} c_{i\sigma}
  \label{eq:E5}
\end{equation}

[E1, N]

\begin{equation}
  0
  \label{eq:E6}
\end{equation}

\textbf{Conservation of $S^z$.} Written the textbook way, $\hat S^z=\sum_{i\sigma}\sigma\,\hat n_{i\sigma}$: the $\sigma$ in the COEFFICIENT position gets the \textbf{value} ($\pm1/2$, an exact rational), the one in the INDEX position gets the \textbf{spelling} (the arrow). Also 0.

$S_z = \sum_{i=1}^{2} \sum_{\sigma} \sigma c^\dagger_{i\sigma} c_{i\sigma}$

\begin{equation}
  S_z = \sum_{i=1}^{2} \sum_{\sigma} \sigma c^\dagger_{i\sigma} c_{i\sigma}
  \label{eq:E7}
\end{equation}

[E1, $S_z$]

\begin{equation}
  0
  \label{eq:E8}
\end{equation}

\textbf{Derive the SU(2) algebra.} Write the ladder operators $\hat S^\pm$ and compute $[\hat S^+,\hat S^-]$. It should be $2\hat S^z$ — and the point is to compute that rather than assume it.

$S_p = \sum_{i=1}^{2} c^\dagger_{i\uparrow} c_{i\downarrow}$

\begin{equation}
  S_p = \sum_{i=1}^{2} c^\dagger_{i\uparrow} c_{i\downarrow}
  \label{eq:E9}
\end{equation}

$S_m = \sum_{i=1}^{2} c^\dagger_{i\downarrow} c_{i\uparrow}$

\begin{equation}
  S_m = \sum_{i=1}^{2} c^\dagger_{i\downarrow} c_{i\uparrow}
  \label{eq:E10}
\end{equation}

[$S_p$, $S_m$]

\begin{equation}
  {c}_{1{\uparrow}}^{\dagger} {c}_{1{\uparrow}} - {c}_{1{\downarrow}}^{\dagger} {c}_{1{\downarrow}} + {c}_{2{\uparrow}}^{\dagger} {c}_{2{\uparrow}} - {c}_{2{\downarrow}}^{\dagger} {c}_{2{\downarrow}}
  \label{eq:E11}
\end{equation}

Write $2\hat S^z$, subtract it, and let the CAS confirm the identity (it gives 0).

$T = 2 \sum_{i=1}^{2} \sum_{\sigma} \sigma c^\dagger_{i\sigma} c_{i\sigma}$

\begin{equation}
  T = 2 \sum_{i=1}^{2} \sum_{\sigma} \sigma c^\dagger_{i\sigma} c_{i\sigma}
  \label{eq:E12}
\end{equation}

subtract E12 from E11

\begin{equation}
  {c}_{1{\uparrow}}^{\dagger} {c}_{1{\uparrow}} - {c}_{1{\downarrow}}^{\dagger} {c}_{1{\downarrow}} + {c}_{2{\uparrow}}^{\dagger} {c}_{2{\uparrow}} - {c}_{2{\downarrow}}^{\dagger} {c}_{2{\downarrow}} - T = 0
  \label{eq:E13}
\end{equation}

\textbf{SU(2) symmetry of $\hat H$}: $[\hat H,\hat S^+]=0$. The $U$ term $\hat n_\uparrow\hat n_\downarrow$ is invariant under flipping the spins, and hopping conserves spin.

[E1, $S_p$]

\begin{equation}
  0
  \label{eq:E14}
\end{equation}

\textbf{Something that is NOT conserved.} The double occupancy $\hat D=\sum_i\hat n_{i\uparrow}\hat n_{i\downarrow}$ is not a constant of the motion. Only the \textbf{hopping} terms survive; $U$ drops out — creating and destroying double occupancy is hopping's doing.

$D = \sum_{i=1}^{2} c^\dagger_{i\uparrow} c_{i\uparrow} c^\dagger_{i\downarrow} c_{i\downarrow}$

\begin{equation}
  D = \sum_{i=1}^{2} c^\dagger_{i\uparrow} c_{i\uparrow} c^\dagger_{i\downarrow} c_{i\downarrow}
  \label{eq:E15}
\end{equation}

[E1, D]

\begin{equation}
  - w {c}_{1{\uparrow}}^{\dagger} {c}_{1{\downarrow}}^{\dagger} {c}_{1{\uparrow}} {c}_{2{\downarrow}} + w {c}_{1{\uparrow}}^{\dagger} {c}_{1{\downarrow}}^{\dagger} {c}_{1{\downarrow}} {c}_{2{\uparrow}} + w {c}_{1{\uparrow}}^{\dagger} {c}_{2{\downarrow}}^{\dagger} {c}_{1{\uparrow}} {c}_{1{\downarrow}} + w {c}_{1{\uparrow}}^{\dagger} {c}_{2{\downarrow}}^{\dagger} {c}_{2{\uparrow}} {c}_{2{\downarrow}} - w {c}_{1{\downarrow}}^{\dagger} {c}_{2{\uparrow}}^{\dagger} {c}_{1{\uparrow}} {c}_{1{\downarrow}} - w {c}_{1{\downarrow}}^{\dagger} {c}_{2{\uparrow}}^{\dagger} {c}_{2{\uparrow}} {c}_{2{\downarrow}} - w {c}_{2{\uparrow}}^{\dagger} {c}_{2{\downarrow}}^{\dagger} {c}_{1{\uparrow}} {c}_{2{\downarrow}} + w {c}_{2{\uparrow}}^{\dagger} {c}_{2{\downarrow}}^{\dagger} {c}_{1{\downarrow}} {c}_{2{\uparrow}}
  \label{eq:E16}
\end{equation}

\textbf{The right-hand side of the equation of motion}, the core of $i\hbar\,\partial_t\hat c_{1\uparrow}=[\hat c_{1\uparrow},\hat H]$: hopping carries it to the neighbouring site, $U$ brings in the opposite-spin number operator.

$[E1, c_{1\uparrow}]$

\begin{equation}
  w {c}_{2{\uparrow}} + U {c}_{1{\downarrow}}^{\dagger} {c}_{1{\uparrow}} {c}_{1{\downarrow}}
  \label{eq:E17}
\end{equation}

\textbf{A nested commutator is read as written}: $[\hat c_{1\uparrow},[\hat c^\dagger_{1\uparrow},\hat c_{1\uparrow}]]=2\hat c_{1\uparrow}$. Which comma separates is decided by bracket depth, so the inner bracket is part of the operand.

$[c_{1\uparrow}, [c^\dagger_{1\uparrow}, c_{1\uparrow}]]$

\begin{equation}
  2 {c}_{1{\uparrow}}
  \label{eq:E18}
\end{equation}

\textbf{Pauli's identity.} Normal-ordering $\hat n^2$ gives $\hat n$ back — two electrons cannot occupy the same state. Another algebraic result, not a transcription.

normal order $c^\dagger_{1\uparrow}$ $c_{1\uparrow}$ $c^\dagger_{1\uparrow}$ $c_{1\uparrow}$

\begin{equation}
  c^\dagger_{1\uparrow} c_{1\uparrow} c^\dagger_{1\uparrow} c_{1\uparrow}
  \label{eq:E19}
\end{equation}

\begin{equation}
  {c}_{1{\uparrow}}^{\dagger} {c}_{1{\uparrow}}
  \label{eq:E20}
\end{equation}

\textbf{Exact time evolution in the atomic limit} (PAIR-REP, docs/34 §2). Switch off the hopping and keep the $U$ term: the $\mathrm{ad}_H$ chain closes on finitely many operators, so $\hat c_{1\uparrow}(t)$ comes out \textbf{with no truncation}.

$H_U = U c^\dagger_{1\uparrow} c_{1\uparrow} c^\dagger_{1\downarrow} c_{1\downarrow}$

\begin{equation}
  H_U = U c^\dagger_{1\uparrow} c_{1\uparrow} c^\dagger_{1\downarrow} c_{1\downarrow}
  \label{eq:E21}
\end{equation}

The coefficient $1-e^{-iUt}$ is a scalar and cannot live in QuantumAlgebra (which decomposes c-numbers into cos/sin) — the substance is the (coefficient, operator) PAIR, and only the display looks like $\hat c_{1\uparrow}(t)$. This is the Hubbard atom's exact result.

evolve $c_{1\uparrow}$ under $H_U$

\begin{equation}
  c_{1\uparrow} + \left(1 - e^{- i U t}\right) c^{\dagger}_{1\downarrow} c_{1\uparrow} c_{1\downarrow}
  \label{eq:E22}
\end{equation}

\end{document}